\documentclass[11pt, a4paper]{article}

% bfw-style.tex — gemeinsame Style-Definitionen für alle Handouts/Aufgabenhefte
% Voraussetzung: Kompilation mit xelatex (wegen fontspec)

\usepackage[ngerman]{babel}
\usepackage{geometry}
\geometry{a4paper, top=2.2cm, bottom=2.2cm, left=2.3cm, right=2.3cm}

\usepackage{fontspec}
% Fallback-Kette: Source Sans Pro -> Calibri -> Segoe UI -> Latin Modern Sans
% (Latin Modern ist immer mit MiKTeX vorhanden)
\IfFontExistsTF{Source Sans Pro}{%
  \setmainfont{Source Sans Pro}[Ligatures=TeX]%
}{\IfFontExistsTF{Calibri}{%
  \setmainfont{Calibri}[Ligatures=TeX]%
}{\IfFontExistsTF{Segoe UI}{%
  \setmainfont{Segoe UI}[Ligatures=TeX]%
}{%
  \setmainfont{Latin Modern Sans}[Ligatures=TeX]%
}}}
\IfFontExistsTF{Source Code Pro}{%
  \setmonofont{Source Code Pro}[Scale=0.92]%
}{\IfFontExistsTF{Consolas}{%
  \setmonofont{Consolas}[Scale=0.92]%
}{%
  \setmonofont{Latin Modern Mono}[Scale=0.92]%
}}

\usepackage{xcolor}
% Farbpalette: hell, freundlich, modern
\definecolor{bfwPrimary}{HTML}{0EA5A4}    % Petrol/Türkis - Primär
\definecolor{bfwPrimaryDark}{HTML}{0F766E} % dunkler Petrol für Headings
\definecolor{bfwAccent}{HTML}{F59E0B}     % Amber - Akzent
\definecolor{bfwSuccess}{HTML}{10B981}    % Grün - Erfolg / Lösung
\definecolor{bfwWarn}{HTML}{F97316}       % Orange - Achtung
\definecolor{bfwInk}{HTML}{0F172A}        % fast Schwarz
\definecolor{bfwInkSoft}{HTML}{475569}    % gedämpftes Grau
\definecolor{bfwBgSoft}{HTML}{F8FAFC}     % off-white
\definecolor{bfwBgTeal}{HTML}{ECFEFF}     % helles Türkis
\definecolor{bfwBgAmber}{HTML}{FEF3C7}    % helles Gelb
\definecolor{bfwBgRose}{HTML}{FFE4E6}     % helles Rosé für Eselsbrücken

\usepackage[colorlinks=true, linkcolor=bfwPrimaryDark, urlcolor=bfwPrimaryDark]{hyperref}
\usepackage{enumitem}
\setlist[itemize]{leftmargin=*, itemsep=2pt, topsep=2pt}
\setlist[enumerate]{leftmargin=*, itemsep=2pt, topsep=2pt}

\usepackage[most]{tcolorbox}
\usepackage{booktabs}
\usepackage{tikz}
\usetikzlibrary{decorations.pathmorphing, positioning, shapes.geometric, arrows.meta}

% Merkkästchen — wichtige Definition / Kernpunkt
\newtcolorbox{merksatz}[1][]{%
  enhanced, breakable,
  colback=bfwBgTeal,
  colframe=bfwPrimary,
  boxrule=0.6pt,
  arc=4pt,
  left=10pt, right=10pt, top=8pt, bottom=8pt,
  fonttitle=\bfseries\color{bfwPrimaryDark},
  title={\faLightbulb\ Merksatz},
  #1
}

% Eselsbrücke — Mnemoniks
\newtcolorbox{eselsbruecke}[1][]{%
  enhanced, breakable,
  colback=bfwBgRose,
  colframe=bfwAccent,
  boxrule=0.6pt,
  arc=4pt,
  left=10pt, right=10pt, top=8pt, bottom=8pt,
  fonttitle=\bfseries\color{bfwWarn},
  title={Eselsbrücke},
  #1
}

% Beispiel-Box
\newtcolorbox{beispiel}[1][]{%
  enhanced, breakable,
  colback=bfwBgSoft,
  colframe=bfwInkSoft,
  boxrule=0.4pt,
  arc=3pt,
  left=10pt, right=10pt, top=8pt, bottom=8pt,
  fonttitle=\bfseries\color{bfwInk},
  title={Beispiel},
  #1
}

% Lösungs-Box (für Aufgabenhefte)
\newtcolorbox{loesung}[1][]{%
  enhanced, breakable,
  colback=bfwBgAmber!40,
  colframe=bfwSuccess,
  boxrule=0.6pt,
  arc=3pt,
  left=10pt, right=10pt, top=8pt, bottom=8pt,
  fonttitle=\bfseries\color{bfwSuccess!70!black},
  title={Lösung},
  #1
}

% Glossar-Eintrag
\newcommand{\glossareintrag}[2]{%
  \par\noindent\textbf{\color{bfwPrimaryDark}#1} \enspace #2\par\smallskip
}

% Section-Stil — etwas farbiger als Standard
\usepackage{titlesec}
\titleformat{\section}
  {\Large\bfseries\color{bfwPrimaryDark}}
  {\thesection}{0.6em}{}
\titleformat{\subsection}
  {\large\bfseries\color{bfwPrimary}}
  {\thesubsection}{0.5em}{}
\titleformat{\subsubsection}
  {\normalsize\bfseries\color{bfwInk}}
  {\thesubsubsection}{0.4em}{}

% TikZ-Stil "skizzenhaft" — etwas weicher als Rough.js, aber stabil
\tikzset{
  roughbox/.style = {
    draw=bfwPrimaryDark, thick, fill=bfwBgTeal,
    rounded corners=4pt, inner sep=6pt
  },
  roughaccent/.style = {
    draw=bfwAccent, thick, fill=bfwBgAmber,
    rounded corners=4pt, inner sep=6pt
  },
  roughline/.style = {
    draw=bfwInkSoft, thick, -{Stealth[length=2.5mm]}
  }
}

% Bullet-Symbole (bewusst kein FontAwesome — vermeidet Font-Installations-Probleme)
\newcommand{\faLightbulb}{$\bullet$}
\newcommand{\faBookmark}{$\bullet$}

% Header/Footer
\usepackage{fancyhdr}
\pagestyle{fancy}
\fancyhf{}
\renewcommand{\headrulewidth}{0pt}
\fancyfoot[L]{\small\color{bfwInkSoft}\thepage}
\fancyfoot[R]{\small\color{bfwInkSoft} BFW M-064 Beschaffung im Büromanagement}


\title{\vspace*{-1.5cm}\Huge\bfseries\color{bfwPrimaryDark}Aufgabenheft Tag~8\\[0.4em]
  \large\color{bfwInkSoft}\textnormal{Besprechungen nachbereiten \& Protokoll --- Übungen im IHK-Klausurformat}}
\author{\textsf{Büroprozesse gestalten und Arbeitsvorgänge organisieren \textperiodcentered{} M-067}\\
\textsf{\small\copyright\ Can Siebert\,\textperiodcentered\,2026}}
\date{Selbstlernmaterial \textperiodcentered{} 08.07.2026}

\begin{document}
\maketitle
\thispagestyle{empty}

\begin{merksatz}[title={\bfseries So nutzen Sie dieses Aufgabenheft}]
Bearbeiten Sie alle Aufgaben zunächst \emph{ohne} in die Lösungen zu schauen. Schreiben
Sie Ihre Antworten in vollständigen Sätzen --- die IHK-Prüfung verlangt formulierte
Begründungen, nicht bloße Stichworte. Beachten Sie die \emph{Operatoren} (nennen,
erläutern, beurteilen, formulieren) --- sie geben den geforderten Antwort-Tiefegrad vor.
Erst danach öffnen Sie den Lösungsteil ab Seite~\pageref{loesungen} und vergleichen.
\end{merksatz}

\tableofcontents
\vspace{1em}
\hrule

\section{Aufgaben}

\subsection*{Aufgabe~1 --- Besprechungen nachbereiten (10 Punkte)}

\noindent\textbf{\color{bfwAccent}YouTube-Vorbereitung:}\enspace \href{https://www.youtube.com/watch?v=1KSQOvp4BJs}{Besprechung nachbereiten --- To-do-Liste} (\url{https://www.youtube.com/watch?v=1KSQOvp4BJs})

\medskip
Sie sind in der Verwaltung der \emph{Lindner \& Partner GmbH} tätig und sollen nach einer
Abteilungsbesprechung die Nachbereitung übernehmen.

\begin{enumerate}[label=(\alph*)]
  \item \textbf{Erläutern} Sie, warum die Nachbereitung einer Besprechung genauso wichtig ist wie ihre Vorbereitung. \hfill (3~P.)
  \item \textbf{Nennen} Sie die drei Angaben, die eine To-do-/Maßnahmenliste je Aufgabe enthalten sollte. \hfill (3~P.)
  \item \textbf{Beschreiben} Sie kurz, wofür der Verteiler eines Protokolls dient. \hfill (2~P.)
  \item \textbf{Beurteilen} Sie die Aussage: \emph{„Wenn die Besprechung vorbei ist, ist die Arbeit getan."} \hfill (2~P.)
\end{enumerate}

\subsection*{Aufgabe~2 --- Funktionen des Protokolls (12 Punkte)}

\noindent\textbf{\color{bfwAccent}YouTube-Vorbereitung:}\enspace \href{https://www.youtube.com/watch?v=SG4tsAG_-hI}{Funktionen eines Protokolls} (\url{https://www.youtube.com/watch?v=SG4tsAG_-hI})

\medskip
Das Wort \emph{Protokoll} bedeutet Niederschrift und meint eine schriftliche Aufzeichnung
über Inhalte, Verlauf und Ergebnisse von Veranstaltungen.

\begin{enumerate}[label=(\alph*)]
  \item \textbf{Definieren} Sie den Begriff „Protokoll". \hfill (3~P.)
  \item \textbf{Nennen} Sie die drei Funktionen, für die Protokolle erstellt werden. \hfill (3~P.)
  \item \textbf{Ordnen} Sie jeder Funktion ein konkretes Beispiel \textbf{zu}. \hfill (3~P.)
  \item \textbf{Erläutern} Sie, warum ein Protokoll als Beweismittel dienen kann. \hfill (3~P.)
\end{enumerate}

\subsection*{Aufgabe~3 --- Protokollarten (12 Punkte)}

\noindent\textbf{\color{bfwAccent}YouTube-Vorbereitung:}\enspace \href{https://www.youtube.com/watch?v=WREpAUH4cs8}{Protokollarten im Überblick} (\url{https://www.youtube.com/watch?v=WREpAUH4cs8})

\medskip
Nach ihrer Ausführlichkeit unterscheidet man Wort-, Verlaufs-, Kurz- und
Beschluss-/Ergebnisprotokoll; daneben gibt es das Gedächtnisprotokoll.

\begin{enumerate}[label=(\alph*)]
  \item \textbf{Erläutern} Sie den Unterschied zwischen \emph{Wortprotokoll} und \emph{Kurzprotokoll}. \hfill (3~P.)
  \item \textbf{Begründen} Sie, welche Protokollart in der kaufmännischen Praxis am häufigsten eingesetzt wird. \hfill (3~P.)
  \item \textbf{Ordnen} Sie den folgenden Anlässen jeweils eine passende Protokollart \textbf{zu}: Gerichtsverhandlung, Projektsitzung, wöchentliche Abteilungsbesprechung. \hfill (3~P.)
  \item \textbf{Erläutern} Sie, was ein Gedächtnisprotokoll von den übrigen Protokollarten unterscheidet. \hfill (3~P.)
\end{enumerate}

\subsection*{Aufgabe~4 --- Protokollrahmen nach DIN 5008 (13 Punkte)}

\noindent\textbf{\color{bfwAccent}YouTube-Vorbereitung:}\enspace \href{https://www.youtube.com/watch?v=qRrRsYK8_PI}{Protokollrahmen nach DIN 5008} (\url{https://www.youtube.com/watch?v=qRrRsYK8_PI})

\medskip
Nach DIN 5008 gliedert sich ein Protokoll in Protokollkopf, Protokolltext und
Protokollfuß.

\begin{enumerate}[label=(\alph*)]
  \item \textbf{Nennen} Sie mindestens fünf Angaben des Protokollkopfes. \hfill (4~P.)
  \item \textbf{Beschreiben} Sie, was der Protokollfuß enthält. \hfill (3~P.)
  \item \textbf{Erläutern} Sie, wie und warum die zeitweise Anwesenheit eines Teilnehmers vermerkt wird. \hfill (3~P.)
  \item \textbf{Erläutern} Sie die drei Kriterien, an denen sich der Protokollführer beim Protokolltext orientiert. \hfill (3~P.)
\end{enumerate}

\subsection*{Aufgabe~5 --- Sprache und indirekte Rede (16 Punkte)}

\noindent\textbf{\color{bfwAccent}YouTube-Vorbereitung:}\enspace \href{https://www.youtube.com/watch?v=fSntdA-DThI}{Indirekte Rede im Protokoll} (\url{https://www.youtube.com/watch?v=fSntdA-DThI})

\medskip
Der Protokolltext wird im Präsens, sachlich und in indirekter Rede formuliert.

\begin{enumerate}[label=(\alph*)]
  \item \textbf{Formulieren} Sie neutral: \emph{„Mit großem Eifer trägt Frau Müller ihr gelungenes Konzept für die Umgestaltung des Raumes vor."} \hfill (3~P.)
  \item \textbf{Wandeln} Sie in indirekte Rede um: \emph{Herr Dreyer: „Ich hole zwei oder drei Angebote für den Betriebsausflug ein."} \hfill (4~P.)
  \item \textbf{Wandeln} Sie in indirekte Rede um: \emph{Frau Bremer: „Die Kunden denken gleich an eine Beschwerde."} (3.\ Person Plural beachten!) \hfill (4~P.)
  \item \textbf{Erläutern} Sie, in welchem Fall der Konjunktiv~II statt des Konjunktivs~I verwendet wird. \hfill (3~P.)
  \item \textbf{Erklären} Sie, warum Anträge und Beschlüsse \emph{wörtlich} ins Protokoll aufgenommen werden. \hfill (2~P.)
\end{enumerate}

\subsection*{Aufgabe~6 --- Protokoll formulieren \& Bedeutung (15 Punkte)}

\noindent\textbf{\color{bfwAccent}YouTube-Vorbereitung:}\enspace \href{https://www.youtube.com/watch?v=aBNW4f6PoS0}{Protokoll genehmigen \& digital führen} (\url{https://www.youtube.com/watch?v=aBNW4f6PoS0})

\medskip
In einer Betriebsratssitzung möchte Frau Käser ein Rundschreiben für die Mitarbeiter
einführen und stellt nach der Diskussion folgenden Antrag: \emph{„Die Mitarbeiter erhalten
monatlich ein Rundschreiben des Betriebsrates mit aktuellen Informationen."} Es wird mit
Handzeichen abgestimmt; alle Betriebsratsmitglieder stimmen für den Antrag.

\begin{enumerate}[label=(\alph*)]
  \item \textbf{Formulieren} Sie den passenden Protokolleintrag (TOP) zu dieser Situation. \hfill (4~P.)
  \item \textbf{Nennen} Sie die zwei Möglichkeiten, ein Protokoll zu genehmigen. \hfill (3~P.)
  \item \textbf{Erläutern} Sie, wodurch die Inhalte eines Protokolls verbindlich werden und welche Protokolle eine Beweisfunktion haben. \hfill (4~P.)
  \item \textbf{Beurteilen} Sie zwei Vorteile und eine Grenze digitaler Protokollvorlagen bzw.\ des kollaborativen Protokollierens. \hfill (4~P.)
\end{enumerate}

\vspace{1em}
\noindent\textbf{Punkte gesamt: 78}

\newpage

\section{Lösungen}\label{loesungen}

\subsection*{Lösung Aufgabe~1}
\begin{loesung}
\textbf{(a)} Ohne Nachbereitung verpuffen die Ergebnisse einer Besprechung: Erst die
Aufgabenverfolgung sorgt dafür, dass beschlossene Maßnahmen tatsächlich umgesetzt werden.
Das Protokoll hält die Ergebnisse fest, informiert Abwesende und macht Entscheidungen
nachvollziehbar --- die investierte Besprechungszeit zahlt sich nur durch konsequente
Nachbereitung aus.

\textbf{(b)} Je Aufgabe gehören in die Maßnahmenliste: \textbf{Was} (Aufgabe/Maßnahme),
\textbf{Wer} (verantwortliche Person) und \textbf{Bis wann} (Termin).

\textbf{(c)} Der Verteiler legt fest, wer das fertige Protokoll erhält --- Teilnehmer,
entschuldigt Abwesende und ggf.\ Vorgesetzte. So sind auch nicht durchgängig Anwesende
über die Ergebnisse informiert.

\textbf{(d)} Die Aussage ist falsch. Mit dem Ende der Besprechung beginnt die
Nachbereitung: Aufgaben werden verfolgt, das Protokoll erstellt und verteilt. Ohne diese
Schritte bleiben Beschlüsse folgenlos.
\end{loesung}

\subsection*{Lösung Aufgabe~2}
\begin{loesung}
\textbf{(a)} Ein Protokoll ist eine \emph{schriftliche Niederschrift} über Inhalte,
Verlauf und Ergebnisse von Veranstaltungen wie Sitzungen, Versammlungen, Besprechungen,
Konferenzen und Verhandlungen.

\textbf{(b)} Die drei Funktionen: \textbf{Information}, \textbf{Beweis- und
Dokumentationsmittel}, \textbf{Organisation}.

\textbf{(c)} Beispiele:
\begin{itemize}[itemsep=0pt]
  \item \emph{Information}: Gedächtnisstütze für Teilnehmer, Information für Abwesende.
  \item \emph{Beweis-/Dokumentationsmittel}: Beschlüsse werden aktenkundig, Entscheidungen u.\,U.\ rechtswirksam.
  \item \emph{Organisation}: Aufgabenverteilung wird festgehalten, Kompetenzen werden festgelegt.
\end{itemize}

\textbf{(d)} Ein Protokoll dokumentiert Beschlüsse und Entscheidungen schriftlich. Mit
\emph{Originalunterschrift} erhält es Beweisfunktion und kann später nachweisen, was mit
welchem Ergebnis verhandelt und beschlossen wurde.
\end{loesung}

\subsection*{Lösung Aufgabe~3}
\begin{loesung}
\textbf{(a)} Das \emph{Wortprotokoll} gibt alle Aussagen \textbf{wörtlich} wieder und ist
sehr umfangreich (Gericht, Parlament). Das \emph{Kurzprotokoll} fasst die wesentlichen
Inhalte \textbf{zusammen}, hält Beschlüsse im Wortlaut mit Abstimmungsergebnissen fest und
verzichtet weitgehend auf Namensnennung.

\textbf{(b)} Am häufigsten ist das \textbf{Kurzprotokoll}: Es stellt Sache und Ergebnis in
den Vordergrund, ist nachvollziehbar und genügt für Abteilungsleiter-, Mitarbeiter- und
Beratungsgespräche --- der typische kaufmännische Alltag.

\textbf{(c)} Zuordnung: Gerichtsverhandlung $\rightarrow$ \emph{Wortprotokoll};
Projektsitzung $\rightarrow$ \emph{Beschluss-/Ergebnisprotokoll}; wöchentliche
Abteilungsbesprechung $\rightarrow$ \emph{Kurzprotokoll}.

\textbf{(d)} Das \emph{Gedächtnisprotokoll} wird erst \textbf{nach} der Besprechung aus dem
Gedächtnis erstellt, wenn währenddessen nicht mitgeschrieben werden konnte; es ist weniger
formal und dient vor allem als persönliche Gedächtnisstütze.
\end{loesung}

\subsection*{Lösung Aufgabe~4}
\begin{loesung}
\textbf{(a)} Angaben des Protokollkopfes (mind.\ fünf): Unternehmen/Abteilung,
Bezeichnung „Protokoll", Besprechungsgruppe/Anlass, Datum und Zeit, Ort, Vorsitz,
Teilnehmer, Protokollführung sowie die Tagesordnung (TOP).

\textbf{(b)} Der Protokollfuß enthält das \emph{Datum der Ausfertigung}, die Namen und
die \emph{Unterschriften} von Leitung (Vorsitz) und Protokollführer; ergänzend können
\emph{Verteiler} und der Hinweis auf \emph{Anlagen} aufgenommen werden.

\textbf{(c)} War ein Teilnehmer nur zeitweise anwesend, wird dies vermerkt (z.\,B.\
„Frau Dr.\ Müller bis 11:20 Uhr, Herr Schütz ab TOP 3"). Wichtig ist das vor allem bei
\emph{Beschlüssen mit rechtlicher Relevanz}, damit nachvollziehbar bleibt, wer bei welcher
Entscheidung anwesend war.

\textbf{(d)} Kriterien für den Protokolltext: \textbf{Wichtigkeit}, \textbf{Vollständigkeit}
und \textbf{wahrheitsgetreue Wiedergabe} der Inhalte und Ergebnisse --- alles ohne Wertung.
\end{loesung}

\subsection*{Lösung Aufgabe~5}
\begin{loesung}
\textbf{(a)} Neutral: „Frau Müller trägt ihr Konzept für die Umgestaltung des Raumes vor."
(Die wertenden Wörter „mit großem Eifer" und „gelungen" entfallen.)

\textbf{(b)} Indirekte Rede (Konjunktiv~I): „Herr Dreyer teilt mit, dass er zwei oder drei
Angebote für den Betriebsausflug \emph{einhole}." (alternativ: „Herr Dreyer erhält den
Auftrag, \dots\ einzuholen.")

\textbf{(c)} In der 3.\ Person Plural ist der Konjunktiv~I („denken") nicht vom Indikativ
zu unterscheiden, daher Konjunktiv~II: „Frau Bremer ergänzt, die Kunden \emph{dächten}
gleich an eine Beschwerde." (Ersatzform: „\dots, die Kunden \emph{würden} gleich an eine
Beschwerde denken.")

\textbf{(d)} Der \emph{Konjunktiv~II} wird gewählt, wenn der Konjunktiv~I formgleich mit
dem Indikativ Präsens ist (außer beim Verb \emph{sein} der Fall bei der 3.\ Person Plural).
In der 3.\ Person Singular ist dagegen der Konjunktiv~I einzusetzen, weil der Konjunktiv~II
dort Zweifel am Wahrheitsgehalt ausdrücken könnte.

\textbf{(e)} Anträge und Beschlüsse werden \emph{wörtlich} aufgenommen, weil ihr genauer
Wortlaut die rechtliche Grundlage bildet --- jede inhaltliche Verschiebung könnte die
Verbindlichkeit oder Auslegung verändern.
\end{loesung}

\subsection*{Lösung Aufgabe~6}
\begin{loesung}
\textbf{(a)} Möglicher Protokolleintrag:
\begin{quote}
\emph{TOP \dots\ Einführung eines Rundschreibens.} Der Betriebsrat beschließt einstimmig:
Die Mitarbeiter erhalten monatlich ein Rundschreiben des Betriebsrates mit aktuellen
Informationen.
\end{quote}
(Der Antragstext fällt weg, es bleibt der Beschlusstext im Wortlaut; Abstimmungsergebnis
„einstimmig" wird genannt.)

\textbf{(b)} Genehmigung auf zwei Arten: (1) Verteilung an die Teilnehmer mit
Widerspruchsfrist --- nach Fristablauf gilt das Protokoll als genehmigt; (2) Vorlesen in
der nächsten Sitzung und Abstimmung über die Genehmigung.

\textbf{(c)} Durch \emph{Unterschrift und Genehmigung} werden Aufträge, Bewilligungen,
Vereinbarungen, Kompetenzen und Beschlüsse \textbf{verbindlich}. Eine \emph{Beweisfunktion}
haben nur Protokolle mit \emph{Originalunterschrift}; sie sind sorgfältig aufzubewahren.

\textbf{(d)} Vorteile: einheitlicher, DIN-gerechter Aufbau über Formularfelder; schnelle
Eingabe schon während der Sitzung; gleichzeitiges Mitschreiben mehrerer Personen;
sofortige Verteilung und Verknüpfung mit Aufgaben-/Terminwerkzeugen. Grenze: Datenschutz
und Zugriffsrechte bzw.\ die fehlende rechtsgültige Unterschrift (ggf.\ qualifizierte
elektronische Signatur erforderlich); auch Ablenkung durch Technik und Abhängigkeit von
Software/Internet sind zu nennen.
\end{loesung}

\end{document}
